% unit: u15
% sources: sources/018-exam-1.txt, sources/113-mit18-06scf11-ex1-pdf-download-file.txt, sources/114-mit18-06scf11-ex1s-pdf-download-file.txt
\section{Exam 1}\label{sec:u15}

The first unit exam checks whether the machinery of Part~I has become reflex: elimination
written with elementary matrices $E_{ij}$, the factorization $A=LU$, the inverse, bases for
the column space and the nullspace, the complete solution of $A\vb{x}=\vb{b}$, and what the
reduced echelon form $R$ says about the columns of $A$. This section collects the four exam
problems as practice problems with complete solutions, preceded by the handful of facts they
draw on and followed by variations to practise on. Try each problem before reading its solution.

\subsection{The facts the exam draws on}

Every problem on the exam is an application of a few statements from
Sections~\ref{sec:u05}--\ref{sec:u11}. For an $m\times n$ matrix $A$ of rank $r$:
\begin{itemize}
  \item $A\vb{x}=\vb{b}$ is solvable exactly when $\vb{b}\in\Col{A}$; it is solvable for
    \emph{every} $\vb{b}\in\R^m$ exactly when $r=m$.
  \item A solution, when one exists, is unique exactly when $\Nul{A}=\{\vb{0}\}$, that is,
    when $r=n$. Otherwise the complete solution is $\vb{x}_p+\vb{x}_n$ with $\vb{x}_n$ ranging
    over the $(n-r)$-dimensional nullspace.
  \item The pivot columns of $A$ (not of $R$) are a basis for $\Col{A}$; the $n-r$ special
    solutions are a basis for $\Nul{A}$.
  \item Row operations do not change the nullspace, so $A$ and $R=\operatorname{rref}(A)$
    satisfy exactly the same column relations.
  \item The inverse of an elimination matrix $E_{ij}$ (subtract $\ell$ times row $j$ from row
    $i$) adds $\ell$ times row $j$ back, and the inverse of a product reverses the order.
\end{itemize}
The first two items combine into the four cases of Table~\ref{tab:u15-cases}, which is the
whole content of Problem~1 below.

\begin{table}[ht]
  \centering
  \begin{tabular}{@{}llll@{}}
    \toprule
    Rank & Shape of $R$ & Solutions of $A\vb{x}=\vb{b}$ & Example \\
    \midrule
    $r=m=n$ & $I$ & exactly one, for every $\vb{b}$ & square invertible \\
    $r=n<m$ & $\begin{bmatrix} I\\ 0\end{bmatrix}$ & zero or one & tall, independent columns \\
    $r=m<n$ & $\begin{bmatrix} I & F\end{bmatrix}$ & infinitely many, for every $\vb{b}$ & wide, independent rows \\
    $r<m,\ r<n$ & $\begin{bmatrix} I & F\\ 0 & 0\end{bmatrix}$ & zero or infinitely many & neither \\
    \bottomrule
  \end{tabular}
  \caption{Rank decides how many solutions $A\vb{x}=\vb{b}$ can have ($F$ is the free part of $R$; columns may need reordering).}
  \label{tab:u15-cases}
\end{table}

\begin{example}[Reading the case from two right-hand sides]\label{ex:u15-case}
Take
\[
  A=\begin{bmatrix}1&2\\2&4\\3&6\end{bmatrix}.
\]
Column~2 is twice column~1, so $r=1$, which is smaller than both $m=3$ and $n=2$: the last
row of the table. Check both possibilities. For $\vb{b}=(1,2,3)$, which is column~1, the
system is solvable with $\vb{x}_p=(1,0)$, and adding any multiple of the special solution
$(-2,1)$ gives infinitely many solutions $\vb{x}=(1-2t,\,t)$. For $\vb{b}=(1,0,0)$,
elimination on the augmented matrix gives
\[
  \left[\begin{array}{cc|c}1&2&1\\2&4&0\\3&6&0\end{array}\right]
  \longrightarrow
  \left[\begin{array}{cc|c}1&2&1\\0&0&-2\\0&0&-3\end{array}\right],
\]
and the equations $0=-2$ and $0=-3$ are impossible, so there is no solution. A single
right-hand side with a unique solution would have forced $r=n$; a single unsolvable one
forces $r<m$.
\end{example}

\subsection{Problem 1: rank from solvability}

\begin{exercise}[Exam problem 1]\label{ex:u15-p1}
An $m\times n$ matrix $A$ has the property that
\[
  A\vb{x}=\begin{bmatrix}1\\1\\1\end{bmatrix}\ \text{has no solution}
  \qquad\text{and}\qquad
  A\vb{x}=\begin{bmatrix}0\\1\\0\end{bmatrix}\ \text{has exactly one solution.}
\]
\begin{enumerate}[(a)]
  \item Give all the information you can about $m$, $n$ and the rank $r$.
  \item Find all solutions of $A\vb{x}=\vb{0}$ and explain.
  \item Write down an example of such a matrix $A$.
\end{enumerate}
\end{exercise}

\begin{solution}
(a) The right-hand sides live in $\R^m$, so $m=3$. Exactly one solution for
$\vb{b}=(0,1,0)$ means there are no free variables: $\Nul{A}=\{\vb{0}\}$ and $r=n$. No
solution for $\vb{b}=(1,1,1)$ means $\Col{A}\neq\R^3$, so $r<m=3$. Together,
$r=n\le 2$: either $n=r=1$ or $n=r=2$ (and $m=3$ in both cases).

(b) Since $r=n$, the columns are independent and $A\vb{x}=\vb{0}$ has only the solution
$\vb{x}=\vb{0}$. (The zero vector is always a solution; independence of the columns says no
other combination of them vanishes.)

(c) Two examples, one for each possible $n$:
\[
  n=1:\ A=\begin{bmatrix}0\\1\\0\end{bmatrix},\qquad
  n=2:\ A=\begin{bmatrix}1&0\\0&1\\0&0\end{bmatrix}.
\]
For $n=1$, $A\vb{x}=(0,1,0)$ is solved only by the scalar $x=1$, and $(1,1,1)$ is not a
multiple of $(0,1,0)$. For $n=2$ the column space is the plane $x_3=0$
(Figure~\ref{fig:u15-p1}), which contains $(0,1,0)=A(0,1)$ but not $(1,1,1)$; the two
columns are independent, so the solution $(0,1)$ is unique.
\end{solution}

\begin{figure}[ht]
  \centering
  \begin{tikzpicture}[x={(-0.55cm,-0.42cm)},y={(1.15cm,0cm)},z={(0cm,1.15cm)},>=Stealth,scale=1.45]
    \fill[blue!8] (-0.3,-0.6,0) -- (2,-0.6,0) -- (2,2.2,0) -- (-0.3,2.2,0) -- cycle;
    \draw[blue!40] (-0.3,-0.6,0) -- (2,-0.6,0) -- (2,2.2,0) -- (-0.3,2.2,0) -- cycle;
    \draw[->,gray] (0,0,0) -- (2.5,0,0) node[below left] {$x_1$};
    \draw[->,gray] (0,0,0) -- (0,2.8,0) node[right] {$x_2$};
    \draw[->,gray] (0,0,0) -- (0,0,1.7) node[above] {$x_3$};
    \draw[->,very thick,blue] (0,0,0) -- (0,1,0) node[below right] {$(0,1,0)\in\Col{A}$};
    \draw[->,very thick,red!75!black] (0,0,0) -- (1,1,1) node[above right] {$(1,1,1)\notin\Col{A}$};
    \draw[dashed,red!75!black] (1,1,1) -- (1,1,0);
    \fill[red!75!black] (1,1,0) circle (1pt);
    \node[blue!70!black,anchor=west] at (2,2.25,0) {$\Col{A}$: the plane $x_3=0$};
  \end{tikzpicture}
  \caption{Problem 1(c) with $n=2$: the column space of $A$ is a plane through the origin in $\R^3$.
  Right-hand sides in the plane are reachable (uniquely, since the columns are independent);
  $(1,1,1)$ sits above the plane and is not.}
  \label{fig:u15-p1}
\end{figure}

\begin{trap}
Each hypothesis controls one inequality: uniqueness for one $\vb{b}$ gives $r=n$ (no free
variables), unsolvability for one $\vb{b}$ gives $r<m$. A frequent slip is to read ``exactly
one solution'' as ``$A$ is square and invertible''; that would make every $\vb{b}$ solvable,
contradicting the first hypothesis.
\end{trap}

\subsection{Problem 2: from elimination matrices to \texorpdfstring{$A^{-1}$}{A inverse}, \texorpdfstring{$A$}{A} and \texorpdfstring{$L$}{L}}

\begin{exercise}[Exam problem 2]\label{ex:u15-p2}
A $3\times 3$ matrix $A$ reduces to $I$ by three row operations, in this order:
$E_{21}$ subtracts $4$ times row~1 from row~2; $E_{31}$ subtracts $3$ times row~1 from
row~3; $E_{23}$ subtracts row~3 from row~2.
\begin{enumerate}[(a)]
  \item Write $A^{-1}$ in terms of the $E$'s, then compute it.
  \item What is the original matrix $A$?
  \item What is the lower triangular factor $L$ in $A=LU$?
\end{enumerate}
\end{exercise}

\begin{solution}
(a) The operations give $E_{23}E_{31}E_{21}A=I$, so $A^{-1}=E_{23}E_{31}E_{21}$, with
\[
  E_{21}=\begin{bmatrix}1&0&0\\-4&1&0\\0&0&1\end{bmatrix},\quad
  E_{31}=\begin{bmatrix}1&0&0\\0&1&0\\-3&0&1\end{bmatrix},\quad
  E_{23}=\begin{bmatrix}1&0&0\\0&1&-1\\0&0&1\end{bmatrix}.
\]
Apply them to $I$ in order. $E_{31}E_{21}$ places $-4$ and $-3$ in column~1; then $E_{23}$
subtracts row~3 from row~2:
\[
  E_{31}E_{21}=\begin{bmatrix}1&0&0\\-4&1&0\\-3&0&1\end{bmatrix},\qquad
  A^{-1}=E_{23}(E_{31}E_{21})=\begin{bmatrix}1&0&0\\-1&1&-1\\-3&0&1\end{bmatrix}.
\]
(b) Invert the product in reverse order, $A=E_{21}^{-1}E_{31}^{-1}E_{23}^{-1}$, where each
inverse just flips the sign of its off-diagonal entry:
\[
  A=\begin{bmatrix}1&0&0\\4&1&0\\0&0&1\end{bmatrix}
    \begin{bmatrix}1&0&0\\0&1&0\\3&0&1\end{bmatrix}
    \begin{bmatrix}1&0&0\\0&1&1\\0&0&1\end{bmatrix}
   =\begin{bmatrix}1&0&0\\4&1&0\\3&0&1\end{bmatrix}
    \begin{bmatrix}1&0&0\\0&1&1\\0&0&1\end{bmatrix}
   =\begin{bmatrix}1&0&0\\4&1&1\\3&0&1\end{bmatrix}.
\]
Check: multiplying this $A$ by the $A^{-1}$ of part (a) gives $I$.

(c) Forward elimination on $A$ uses only $E_{21}$ and $E_{31}$: subtracting $4$ (row~1) from
row~2 and $3$ (row~1) from row~3 leaves a matrix that is already upper triangular. The step $E_{23}$ is an \emph{upward} step (it clears the entry above the
last pivot), so it is not part of $A=LU$. The multipliers $4$ and $3$ go straight into $L$:
\[
  L=E_{21}^{-1}E_{31}^{-1}=\begin{bmatrix}1&0&0\\4&1&0\\3&0&1\end{bmatrix},\qquad
  U=E_{23}^{-1}=\begin{bmatrix}1&0&0\\0&1&1\\0&0&1\end{bmatrix},
\]
and indeed the factorization in (b) already reads $A=LU$.
\end{solution}

\begin{trap}
$A^{-1}$ is the product of the $E$'s in the order they were \emph{applied}
($E_{23}E_{31}E_{21}$, last step on the left); $A$ is the product of the inverses in the
\emph{opposite} order. Writing $A^{-1}=E_{21}E_{31}E_{23}$ gives a different matrix here,
because $E_{23}$ does not commute with $E_{31}$.
\end{trap}

\subsection{Problem 3: bases and the complete solution for a matrix with a parameter}

\begin{exercise}[Exam problem 3]\label{ex:u15-p3}
This $3\times 4$ matrix depends on $c$:
\[
  A=\begin{bmatrix}1&1&2&4\\3&c&2&8\\0&0&2&2\end{bmatrix}.
\]
\begin{enumerate}[(a)]
  \item For each $c$ find a basis for the column space of $A$.
  \item For each $c$ find a basis for the nullspace of $A$.
  \item For each $c$ find the complete solution of $A\vb{x}=(1,c,0)$.
\end{enumerate}
\end{exercise}

\begin{solution}
Subtract $3$ (row~1) from row~2:
\[
  \begin{bmatrix}1&1&2&4\\0&c-3&-4&-4\\0&0&2&2\end{bmatrix}.
\]
Everything depends on whether the entry $c-3$ is a pivot.

\emph{Case $c\neq 3$.} The pivots are in columns 1, 2, 3, so $r=3$. Dividing and clearing
upward gives
\[
  R=\begin{bmatrix}1&0&0&2\\0&1&0&0\\0&0&1&1\end{bmatrix}.
\]
(a) The pivot columns of $A$, namely $(1,3,0)$, $(1,c,0)$, $(2,2,2)$, are a basis for
$\Col{A}$. Since $r=m=3$ this column space is all of $\R^3$, so the standard basis would do
equally well.
(b) One free variable, $x_4$. Setting $x_4=1$ in $R\vb{x}=\vb{0}$ gives $x_1=-2$, $x_2=0$,
$x_3=-1$; the basis is $\vb{s}=(-2,0,-1,1)$.

\emph{Case $c=3$.} Row~2 becomes $(0,0,-4,-4)$; adding half of it to row~3 kills row~3, so
$r=2$ with pivots in columns 1 and 3:
\[
  R=\begin{bmatrix}1&1&0&2\\0&0&1&1\\0&0&0&0\end{bmatrix}.
\]
(a) Basis for $\Col{A}$: columns 1 and 3 of $A$, $(1,3,0)$ and $(2,2,2)$.
(b) Free variables $x_2$ and $x_4$. The special solutions are
$\vb{s}_1=(-1,1,0,0)$ (from $x_2=1,x_4=0$) and $\vb{s}_2=(-2,0,-1,1)$ (from $x_2=0,x_4=1$).

(c) The right-hand side $(1,c,0)$ is exactly column~2 of $A$, so $\vb{x}_p=(0,1,0,0)$ is a
particular solution for every $c$. Adding the nullspace:
\[
  c\neq 3:\ \ \vb{x}=\begin{bmatrix}0\\1\\0\\0\end{bmatrix}+x_4\begin{bmatrix}-2\\0\\-1\\1\end{bmatrix};
  \qquad
  c=3:\ \ \vb{x}=\begin{bmatrix}0\\1\\0\\0\end{bmatrix}+x_2\begin{bmatrix}-1\\1\\0\\0\end{bmatrix}
   +x_4\begin{bmatrix}-2\\0\\-1\\1\end{bmatrix}.
\]
When $c=3$ the vector $(1,0,0,0)$ is another particular solution (column~1 equals
$(1,3,0)$ too); it differs from $(0,1,0,0)$ by the null vector $\vb{s}_1$, so it describes
the same solution set.
\end{solution}

\begin{trap}
The basis for $\Col{A}$ consists of columns of $A$, not of $R$. For $c=3$ the pivot columns
of $R$ are $(1,0,0)$ and $(0,1,0)$, which span the plane $x_3=0$. But column~3 of $A$ is
$(2,2,2)$, which is not in that plane, so those two vectors are not a basis for $\Col{A}$.
Row operations preserve the \emph{relations} among the columns, not the column space.
\end{trap}

\subsection{Problem 4: what \texorpdfstring{$\operatorname{rref}(A)$}{rref(A)} reveals}

\begin{exercise}[Exam problem 4]\label{ex:u15-p4}
\begin{enumerate}[(a)]
  \item If $A$ is $3\times 5$, what do you know about its nullspace?
  \item Row operations on $A$ lead to
    \[
      R=\operatorname{rref}(A)=\begin{bmatrix}1&4&0&0&0\\0&0&0&1&0\\0&0&0&0&1\end{bmatrix}.
    \]
    Write down everything you know about the columns of $A$.
  \item In the vector space $M$ of all $3\times 3$ matrices, which subspace $S$ is spanned by
    all possible reduced row echelon forms $R$?
\end{enumerate}
\end{exercise}

\begin{solution}
(a) $\Nul{A}$ is a subspace of $\R^5$ of dimension $5-r$. Since $r\le 3$, the dimension is
at least $2$: the nullspace contains at least two independent vectors, and every
$3\times 5$ matrix has nonzero solutions to $A\vb{x}=\vb{0}$.

(b) Write $\vb{a}_1,\dots,\vb{a}_5$ for the columns of $A$. The pivots of $R$ are in columns
1, 4, 5, so $r=3$ and:
\begin{itemize}
  \item $\vb{a}_1,\vb{a}_4,\vb{a}_5$ are independent and form a basis for
    $\Col{A}=\R^3$ (three pivots in three rows);
  \item $\vb{a}_2=4\vb{a}_1$, because $(-4,1,0,0,0)\in\Nul{R}=\Nul{A}$ means
    $-4\vb{a}_1+\vb{a}_2=\vb{0}$;
  \item $\vb{a}_3=\vb{0}$, because $(0,0,1,0,0)\in\Nul{A}$.
\end{itemize}
Nothing else is determined: any three independent vectors in $\R^3$ can serve as
$\vb{a}_1,\vb{a}_4,\vb{a}_5$.

(c) In an echelon matrix the pivot of row $i$ sits in a column $p_i$ with
$p_1<p_2<\cdots$, so $p_i\ge i$ and everything to the left of a pivot is zero; zero rows are
at the bottom. Hence every $R$ is \emph{upper triangular}, and so is every combination of
$R$'s. Conversely, the six matrices $E_{ij}$ ($i\le j$) with a single $1$ in position
$(i,j)$ are combinations of echelon forms:
$E_{11}$, $E_{12}$, $E_{13}$ are themselves echelon forms; $E_{22}$ is the echelon form
$\left[\begin{smallmatrix}1&0&0\\0&1&0\\0&0&0\end{smallmatrix}\right]$ minus $E_{11}$;
$E_{23}$ is $\left[\begin{smallmatrix}0&1&0\\0&0&1\\0&0&0\end{smallmatrix}\right]$ minus
$E_{12}$; $E_{33}$ is $I$ minus
$\left[\begin{smallmatrix}1&0&0\\0&1&0\\0&0&0\end{smallmatrix}\right]$.
So $S$ is the space of all $3\times 3$ upper triangular matrices, of dimension $6$.
\end{solution}

\subsection{Variations for practice}

\begin{exercise}\label{ex:u15-v1}
Suppose $A\vb{x}=\vb{b}$ has a solution for every $\vb{b}\in\R^2$, and $A\vb{x}=(1,0)$ has
infinitely many solutions. What can you say about $m$, $n$, $r$? Give an example.
\end{exercise}

\begin{solution}
$\vb{b}\in\R^2$ gives $m=2$; solvable for every $\vb{b}$ gives $r=m=2$; infinitely many
solutions means free variables, so $n>r=2$. Thus $m=r=2<n$. Example:
$A=\begin{bmatrix}1&0&0\\0&1&0\end{bmatrix}$, where $A\vb{x}=(1,0)$ is solved by
$(1,0,t)$ for every $t$.
\end{solution}

\begin{exercise}\label{ex:u15-v2}
Elimination on a $3\times 3$ matrix $A$ uses $E_{21}$ (subtract $2$ times row~1 from
row~2) and then $E_{32}$ (subtract $5$ times row~2 from row~3), reaching
$U=\begin{bmatrix}1&1&1\\0&1&1\\0&0&1\end{bmatrix}$. Find $L$ and $A$.
\end{exercise}

\begin{solution}
In the order $E_{21}$, $E_{32}$ the multipliers go directly into $L$:
$L=E_{21}^{-1}E_{32}^{-1}=\begin{bmatrix}1&0&0\\2&1&0\\0&5&1\end{bmatrix}$. Then
\[
  A=LU=\begin{bmatrix}1&0&0\\2&1&0\\0&5&1\end{bmatrix}
       \begin{bmatrix}1&1&1\\0&1&1\\0&0&1\end{bmatrix}
      =\begin{bmatrix}1&1&1\\2&3&3\\0&5&6\end{bmatrix}.
\]
Check: subtracting $2$ (row~1) from row~2 gives $(0,1,1)$, then subtracting $5$ times that
from row~3 gives $(0,0,1)$.
\end{solution}
